\addcontentsline{toc}{section}{CAPÍTULO XII}
\chapter{APRENDIZAJES Y LECCIONES}

\section{Técnicos}

Los modelos generalistas de frontera (\textit{Claude Sonnet}, \textit{Claude Opus}) superaron al modelo especialista médico (MedGemma 27B) en una tarea que, en principio, debería favorecer al especialista. MedGemma está entrenado con corpus médicos y entiende terminología clínica. Sin embargo, la tarea de auditoría de liquidaciones no es puramente médica: requiere razonamiento sobre datos tabulares heterogéneos (códigos, descripciones, costos), detección de patrones administrativos (reversiones, duplicados) y evaluación de proporcionalidad económica. Estas capacidades dependen más de la escala y la versatilidad del modelo que de su especialización en el dominio clínico. La ventana de contexto de 8,192 \textit{tokens} de MedGemma fue una limitación adicional que obligó a truncar la información de entrada, perdiendo visibilidad sobre los ítems de cola larga donde frecuentemente se esconden los sobrecostos.

Dos hallazgos operativos complementan esta lección. La desviación estándar, como métrica de discriminación, resultó ser el indicador más predictivo de utilidad: si el equipo hubiese seleccionado por media de puntaje o de número de hallazgos, habría elegido Opus, un modelo que discrimina poco y no genera clasificaciones \textit{fast-track}. Y el pre-filtrado de Fase 1 aportó un valor imprevisto pero contra-intuitivo: al inyectar el contexto de cobertura contractual en el \textit{prompt} del LLM, no redujo la tasa de sobre-escalación sino que la aumentó en 7 pp (de 72.3 \% a 65.3 \% de fast-track, según el estudio de ablación de §10.2). El valor de Fase 1 no está en reducir falsos positivos sino en mejorar la calidad de la discriminación: los 7 \textit{claims} adicionales que se escalan son precisamente aquellos con discrepancias contractuales que el LLM, sin contexto de cobertura, habría ignorado. La sinergia entre el motor de reglas y el LLM no es aditiva (ambos reducen riesgo) sino complementaria (cada uno detecta un tipo distinto de anomalía).

\section{Metodológicos}

El \textit{Golden Set} de 17 liquidaciones resultó más útil de lo esperado para un conjunto tan pequeño. Su función no es estadística (17 casos no prueban nada en sentido inferencial), sino de calibración y regresión. Lo que le da credibilidad es que las clasificaciones esperadas fueron definidas por un auditor médico, no por el equipo técnico: el \textit{Golden Set} refleja criterio clínico experto, no una expectativa autorreferencial. Cada vez que se modificó un umbral, se ajustó un \textit{prompt} o se cambió la lógica de \textit{routing}, el \textit{Golden Set} respondió a la pregunta operativa inmediata: “¿este cambio rompe algún caso que antes funcionaba?” Los 59 \textit{tests} del \textit{Golden Set}   (test golden set.py), 51 parametrizados sobre los 17 \textit{claims}, más 6 de simulación de decisión y 2 de proyección de tasa, se ejecutan en menos de 3 segundos y cubren \textit{routing}, elegibilidad, cobertura y la decisión final. Sin este mecanismo, la calibración de umbrales habría sido un proceso de ensayo y error sin retroalimentación inmediata.

La estrategia escalonada (Sonnet para \textit{screening}, Opus para análisis profundo) constituye un patrón reutilizable más allá de este proyecto. El principio subyacente (usar un modelo rápido y barato para filtrar, y uno lento y caro para profundizar) aplica a cualquier flujo de trabajo donde la distribución de complejidad sea asimétrica (muchos casos simples, pocos complejos). Los criterios de escalación (nexuscare/pipeline/escalation.py) encapsulan la lógica de forma parametrizable: umbrales de \textit{score}, tipos de hallazgo, costo total del caso. Esto permite ajustar la agresividad de la escalación según el contexto operativo sin modificar la arquitectura.

\section{De negocio}

La heterogeneidad de datos que se describe en la sección 2.3 como un problema teórico resultó ser una barrera práctica concreta durante el desarrollo. La extracción de reglas de plan desde los contratos PDF de MediVida consumió más tiempo del previsto porque los contratos utilizan lenguaje ambiguo (“se cubrirán los medicamentos indicados excepto los expresamente excluidos”), referencias cruzadas entre secciones y tablas con formatos inconsistentes. El proceso semiautomático (LLM extrae, humano valida) redujo el esfuerzo, pero cada nuevo plan requiere entre 4 y 8 horas de configuración y validación. En un escenario con decenas de planes activos, como operan las EPS grandes, esto implica un costo operativo recurrente que debe dimensionarse en el modelo de negocio.

La brecha entre la viabilidad técnica y la adopción operativa se hizo evidente durante las conversaciones informales con auditores. El sistema puede procesar 101 liquidaciones en menos de 2 horas; la pregunta operativa no es “¿funciona?” sino “¿quién firma la decisión de clasificar como fast-track ?” En el contexto regulatorio peruano, una liquidación clasificada como fast-track que resulte en una glosa de SUSALUD recae sobre la persona que autorizó la aprobación, no sobre el sistema. Hasta que la normativa se adapte al uso de IA como herramienta de soporte, o hasta que la experiencia acumulada genere confianza institucional, la fast-track operará como una pre-aprobación que el auditor ratifica con un clic, no como una decisión final sin intervención humana.
